\documentclass[10pt, letterpaper]{article}

%------------------------------------------------------------------------------
%  This file contains Zhaoxiang (Aaron) Feng's curriculum vitae rendered using
%  the RenderCV style.
%------------------------------------------------------------------------------

% Packages:
\usepackage[
    ignoreheadfoot, % set margins without considering header and footer
    top=2 cm,      % separation between body and page edge from the top
    bottom=2 cm,   % separation between body and page edge from the bottom
    left=2 cm,     % separation between body and page edge from the left
    right=2 cm,    % separation between body and page edge from the right
    footskip=1.0 cm % separation between body and footer
    % showframe % uncomment for debugging layout
]{geometry}
\usepackage{titlesec}        % customize section titles
\usepackage{tabularx}        % fixed width tables
\usepackage{array}           % required by tabularx
\usepackage[dvipsnames]{xcolor} % color support
\definecolor{primaryColor}{RGB}{0, 79, 144} % RenderCV primary colour
\usepackage{enumitem}        % list customisation
\usepackage{fontawesome5}    % icon fonts
\usepackage{amsmath}         % mathematics support
\usepackage[
    pdftitle={Zhaoxiang (Aaron) Feng's CV},
    pdfauthor={Zhaoxiang Feng},
    pdfcreator={LaTeX with RenderCV},
    colorlinks=true,
    urlcolor=primaryColor
]{hyperref}                  % hyperlinks and PDF metadata
\usepackage[pscoord]{eso-pic}    % floating text on the page
\usepackage{calc}                % length calculations
\usepackage{bookmark}            % PDF bookmarks
\usepackage{lastpage}            % reference the number of pages
\usepackage{changepage}          % adjust width for entries
\usepackage{paracol}             % multi‑column support
\usepackage{ifthen}              % conditional logic
\usepackage{needspace}           % avoid page breaks after titles
\usepackage{iftex}               % engine detection

% Ensure that the generated PDF is machine‑readable and ATS‑parsable.
\ifPDFTeX
    \input{glyphtounicode}
    \pdfgentounicode=1
    \usepackage[utf8]{inputenc}
    \usepackage[T1]{fontenc}
    \usepackage{lmodern}
\fi

%------------------------------------------------------------------------------
%  Layout and formatting settings
%------------------------------------------------------------------------------
\AtBeginEnvironment{adjustwidth}{\partopsep0pt} % remove extra space before adjustwidth
\pagestyle{empty}            % suppress headers and footers
\setcounter{secnumdepth}{0}  % remove section numbering
\setlength{\parindent}{0pt} % no paragraph indentation
\setlength{\topskip}{0pt}   % no top skip
\setlength{\columnsep}{0cm} % spacing between paracol columns

% Customise the section headings to include a horizontal rule beneath.
\titleformat{\section}{\needspace{4\baselineskip}\bfseries\large}{}{0pt}{}[\vspace{1pt}\titlerule]
\titlespacing{\section}{-1pt}{0.3 cm}{0.2 cm}

% Define bullet appearance for the highlights environment.
\renewcommand\labelitemi{$\circ$}
\newenvironment{highlights}{
    \begin{itemize}[
        topsep=0.10 cm,
        parsep=0.10 cm,
        partopsep=0pt,
        itemsep=0pt,
        leftmargin=0.4 cm + 10pt
    ]
}{\end{itemize}}
\newenvironment{highlightsforbulletentries}{
    \begin{itemize}[
        topsep=0.10 cm,
        parsep=0.10 cm,
        partopsep=0pt,
        itemsep=0pt,
        leftmargin=10pt
    ]
}{\end{itemize}}

% Single and two column entry environments.
\newenvironment{onecolentry}{
    \begin{adjustwidth}{0.2 cm + 0.00001 cm}{0.2 cm + 0.00001 cm}
}{\end{adjustwidth}}

\newenvironment{twocolentry}[2][]{
    \onecolentry
    \def\secondColumn{#2}
    \setcolumnwidth{\fill, 4.5 cm}
    \begin{paracol}{2}
}{
    \switchcolumn \raggedleft \secondColumn
    \end{paracol}
    \endonecolentry
}

% Header environment for the CV title and personal details.
\newenvironment{header}{
    \setlength{\topsep}{0pt}\par\kern\topsep\centering\linespread{1.5}
}{
    \par\kern\topsep
}

% Fix for the missing \AND command (renders as a separator pipe)
\newcommand{\AND}{\unskip\enspace\textbar\enspace}

% Save the original \href command and redefine \href to omit arrow icons.
\let\hrefWithoutArrow\href
\renewcommand{\href}[2]{\hrefWithoutArrow{#1}{\ifthenelse{\equal{#2}{}}{ }{#2}}}

\begin{document}

    %--------------------------------------------------------------------------
    %  Header
    %--------------------------------------------------------------------------
    \begin{header}
        \textbf{\fontsize{24 pt}{24 pt}\selectfont Zhaoxiang (Aaron) Feng}

        \vspace{0.3 cm}

        \normalsize
        % The \AND macro inside the header environment draws separators.
        \mbox{{\color{black}\footnotesize\faMapMarker*}\hspace*{0.13cm}San Diego, CA}%
        \kern 0.25 cm% spacing between columns
        \AND%
        \kern 0.25 cm%
        \mbox{\hrefWithoutArrow{mailto:zhf004@ucsd.edu}{\color{black}{\footnotesize\faEnvelope[regular]}\hspace*{0.13cm}zhf004@ucsd.edu}}%
        \kern 0.25 cm%
        \AND%
        \kern 0.25 cm%
        \mbox{\hrefWithoutArrow{tel:+13363377139}{\color{black}{\footnotesize\faPhone*}\hspace*{0.13cm}(336)\,337\,-7139}}%
        \kern 0.25 cm%
        \AND%
        \kern 0.25 cm%
        \mbox{\hrefWithoutArrow{https://github.com/aaronzhfeng}{\color{black}{\footnotesize\faGithub}\hspace*{0.13cm}aaronzhfeng}}%
        \kern 0.25 cm%
        \AND%
        \kern 0.25 cm%        
        \mbox{\hrefWithoutArrow{https://aaronzhfeng.github.io/}{\color{black}{\footnotesize\faGlobe}\hspace*{0.13cm}aaronzhfeng.github.io}}%

    \end{header}

    \vspace{0.3 cm - 0.3 cm}

    %--------------------------------------------------------------------------
    %  Education Section
    %--------------------------------------------------------------------------
    \section{Education}

    % UC San Diego education entry
    \begin{twocolentry}{\textit{Sep. 2022 -- Apr. 2026 (expected)}}
        \textbf{University of California San Diego (UCSD)}

        \textit{B.S. in Data Science \& B.S. in Probability \& Statistics}
    \end{twocolentry}

    \vspace{0.10 cm}
    \begin{onecolentry}
        \begin{highlights}
            \item GPA: 3.83 (Major GPAs: 3.95 in Data Science, 3.93 in Statistics)
            \item Honors: Provost Honors (multiple quarters)
            \item Graduate-level Coursework: Advanced Time Series Analysis, Machine Learning, Applied Statistics, Privacy-sensitive Data Systems
        \end{highlights}
    \end{onecolentry}

    %--------------------------------------------------------------------------
    %  Research Interests Section
    %--------------------------------------------------------------------------
    \section{Research Interests}

    \begin{onecolentry}
        I am interested in building AI agents with persistent \textbf{memory} that enables test-time continual learning: systems that accumulate and reuse concepts, self-verify, and improve after deployment. My focus areas include \textbf{memory} architectures, calibrated uncertainty, and efficient inference.
    \end{onecolentry}

    %--------------------------------------------------------------------------
    %  Publication Section
    %--------------------------------------------------------------------------
    \section{Publications}
    
    \begin{onecolentry}
        [1] Matthew Ho, Chen Si, \textbf{Zhaoxiang Feng}, Fangxu Yu, Yichi Yang, Zhijian Liu, Zhiting Hu, Lianhui Qin. ``ArcMemo: Abstract Reasoning Composition with Lifelong LLM Memory.''\\
        \textit{\href{https://arcprize.org/blog/arc-prize-2025-results-analysis}{\textbf{Runner Up}}, ARC Prize 2025 Paper Awards --- Top 8 of 90 paper submissions.} \hfill \textit{[\href{https://arxiv.org/abs/2509.04439}{paper}]} \textit{[\href{https://github.com/matt-seb-ho/arc_memo}{code}]}
    \end{onecolentry}

    \vspace{0.20 cm}

    \begin{onecolentry}
        [2] \textbf{Zhaoxiang Feng}, David Scott Lewis, Enrique Zueco. ``LabMemo: Concept-Level Memory for Autonomous Scientific Discovery.''\\
        \textit{Accepted at IEEE IROS 2025 Workshop on Embodied AI and Robotics for Future Scientific Discovery (AIR4S) (non-archival)}, 2025. \hfill \textit{[\href{https://github.com/aaronzhfeng/lab_memo/blob/main/paper/LabMemo_Paper.pdf}{paper}]} \textit{[\href{https://github.com/aaronzhfeng/lab_memo}{code}]}
    \end{onecolentry}

    \vspace{0.20 cm}

    \begin{onecolentry}
        [3] \textbf{Zhaoxiang Feng}, David Scott Lewis. ``SOKRATES: Distilling Symbolic Knowledge into Option-Level Reasoning via Solver-Guided Preference Optimization.''\\
        \textit{Accepted at AAAI 2026 Bridge Program on Logic \& AI: Logical and Symbolic Reasoning in Language Models (LMReasoning)}, 2026. \hfill \textit{[\href{https://github.com/aaronzhfeng/sokrates-oak/blob/main/paper/sokrates.pdf}{paper}]} \textit{[\href{https://github.com/aaronzhfeng/sokrates-oak}{code}]}
    \end{onecolentry}

    \vspace{0.20 cm}

    \begin{onecolentry}
        [4] \textbf{Zhaoxiang Feng}, Mingyang Yao, Charlie Sun, David Scott Lewis. ``Which Memory Operation Drives Recovery? A Factorial Study of Retrieve, Write, and Manage Adaptation under Domain Shift.''\\
        \textit{Accepted at ICLR 2026 Workshop on Lifelong Agents: Learning, Aligning, Evolving (LLA)}, 2026. \hfill \textit{[\href{https://openreview.net/pdf?id=81CEtTu1oA}{paper}]}
    \end{onecolentry}

    \vspace{0 cm}
    \begin{onecolentry}\hspace{0.3cm}\rule{1.2cm}{0.3pt}\end{onecolentry}
    \vspace{0.10 cm}

    \begin{onecolentry}
        [5] Jinzhou Tang, Yufan Zhou, Zixuan Wang, \textbf{Zhaoxiang Feng}, Xinle Yu, Steven Ngo, Zhengding Hu, Luoshang Pan, Lianhui Qin, Yufei Ding, Tianmin Shu, Jingbo Shang, Zhiting Hu, Zhen Wang. ``$S^3$-Sim: Simulating Humans for Personalized Language Modeling.''\\
        \textit{Under review at International Conference on Machine Learning (ICML)}, 2026.
    \end{onecolentry}

    \vspace{0.20 cm}

    \begin{onecolentry}
        [6] \textbf{Zhaoxiang Feng}, Zhongyan Luo, Mingyang Yao, Charlie Sun. ``Explanation-Consistency Graphs: Neighborhood Surprise in Explanation Space for Training Data Debugging.''\\
        \textit{Under review at ACL Rolling Review (ARR)}, 2026.
    \end{onecolentry}

    \vspace{0.20 cm}

    \begin{onecolentry}
        [7] \textbf{Zhaoxiang Feng}, Zhongyan Luo, Mingyang Yao, Charlie Sun, Yuwen Zhang, David Scott Lewis. ``EP-Prior: Interpretable ECG Representations via Electrophysiology Constraints.''\\
        \textit{Under review at IJCAI-ECAI 2026 AI and Health Track}, 2026.
    \end{onecolentry}

    \vspace{0.20 cm}

    \begin{onecolentry}
        [8] \textbf{Zhaoxiang Feng}, Mingyang Yao, Shuheng Cao, Young Su Ko, Wei Wang. ``TrustPPI: Deformation Stability as a Model-Agnostic Trust Signal for Protein-Protein Interaction Prediction.''\\
        \textit{Under review at KDD 2026 AI4Sciences Track}, 2026. \hfill \textit{[\href{https://github.com/aaronzhfeng/trust_ppi}{code}]}
    \end{onecolentry}

    %--------------------------------------------------------------------------
    %  Research Experience Section
    %--------------------------------------------------------------------------
    \section{Research Experience}

    % Q-Lab research position
    \begin{twocolentry}{\textit{Jan. 2025 -- Present}}
        \textbf{Research Assistant | Q-Lab, UC San Diego}\\
        \textit{Advisors: Prof. Lianhui Qin and Matthew Ho}
    \end{twocolentry}
    \vspace{0.10 cm}
    \begin{onecolentry}
        \begin{highlights}
            \item Led design of ArcMemo's program-synthesis-style memory ontology: many-to-one puzzle-to-feature mappings and curated concept parameterizations for concept-level reasoning.
            \item Engineered a reasoning-based retrieval mechanism (System-2 exploration) to resolve embedding failures, achieving a 7.5\% relative gain on ARC-AGI-1 (59.33\% official score).
            \item Built a concept dataset generation pipeline transforming hand-written concepts into validated helper puzzles via multi-stage LLM generation, code synthesis, and automated testing.
            \item Extended to AIME math problems via a metacognitive self-assessment pipeline, improving accuracy by 9.3\% through self-reflective memory usage.
        \end{highlights}
    \end{onecolentry}
    \vspace{0.20 cm}

    % Wang Lab Student Researcher
    \begin{twocolentry}{\textit{Jun. 2025 -- Present}}
        \textbf{Research Assistant | Wang Lab, UC San Diego}\\
        \textit{Advisors: Prof. Wei Wang and Young Su Ko}
    \end{twocolentry}
    \vspace{0.10 cm}
    \begin{onecolentry}
        \begin{highlights}
            \item Developing TrustPPI: a model-agnostic trust signal for PPI prediction via deformation stability — injecting Gaussian noise ($\sigma$=0.10, K=10) into intermediate embeddings and combining flip-rate and variance signals; validated across 4 architectures (CNN, Transformer, ESM-2, E(3)-GNN) and 11 model$\times$species combinations.
            \item Achieved 0.66--0.83 AUROC for error detection across all models, with trust-guided selection discovering true interactions 2--3.5$\times$ faster than confidence baselines and +2--10\% accuracy gain at 80\% coverage.
            \item Architected a heterogeneous Mixture-of-Experts system for chemical reaction prediction on USPTO (1M+ reactions) using directed message-passing GNN encoders with shortest-path positional encodings and learned routing mechanisms across four specialized expert models.
        \end{highlights}
    \end{onecolentry}
    \vspace{0.20 cm}

    % HDSI Capstone — Hao Zhang
    \begin{twocolentry}{\textit{Sep. 2025 -- Present}}
        \textbf{Student Researcher | Halıcıo\u{g}lu Data Science Institute, UC San Diego}\\
        \textit{Advisor: Prof. Hao Zhang (Senior Capstone)}
    \end{twocolentry}
    \vspace{0.10 cm}
    \begin{onecolentry}
        \begin{highlights}
            \item Built an end-to-end distributed training stack for a 1.8B-parameter language model on 8 NVIDIA B200 GPUs, implementing pipeline parallelism and scaling analysis.
            \item Ported DFlash (diffusion-based speculative decoding) to JAX/TPU, achieving 3.02$\times$ average speedup (4.93$\times$ on math, 9 benchmarks); exploited TPU MXU 128$\times$128 tile structure for flat verification cost up to K=128, lifting the GPU-motivated K=16 constraint; integrated into vLLM with 2.31$\times$ end-to-end speedup; PR-ready for vllm-project/tpu-inference.
        \end{highlights}
    \end{onecolentry}
    \vspace{0.20 cm}

    % HDSI Capstone — Yian Ma
    \begin{twocolentry}{\textit{Sep. 2025 -- Present}}
        \textbf{Student Researcher | Halıcıo\u{g}lu Data Science Institute, UC San Diego}\\
        \textit{Advisor: Prof. Yian Ma (Senior Capstone)}
    \end{twocolentry}
    \vspace{0.10 cm}
    \begin{onecolentry}
        \begin{highlights}
            \item Designed MI-based uncertainty signals for epistemic uncertainty detection and risk-controlled stopping in LLM reasoning.
            \item Developing an uncertainty-guided interactive reasoning framework using self-revision mutual information for ask-vs-answer stopping; designed Evidence Sensitivity and Hypothesis Discrimination probes with Clopper-Pearson risk-controlled calibration, improving AR-Bench Detective Cases from 36\% to 45\%.
        \end{highlights}
    \end{onecolentry}
    \vspace{0.20 cm}

    % USIPC Research Assistant
    \begin{twocolentry}{\textit{Sep. 2023 -- Jun. 2024}}
        \textbf{Research Assistant | U.S. Immigration Policy Center, UC San Diego}\\
        \textit{Advisors: Prof. Tom K. Wong and Dr. Gabriel De Roche}
    \end{twocolentry}
    \vspace{0.10 cm}
    \begin{onecolentry}
        \begin{highlights}
            \item Built ARIMA and exponential smoothing forecasting models on 20+ years of naturalization data for 2024 election policy briefs; developed automated data pipelines and Tableau dashboards for non-technical stakeholders.
        \end{highlights}
    \end{onecolentry}

    %--------------------------------------------------------------------------
    %  Teaching & Service Section
    %--------------------------------------------------------------------------
    \section{Teaching \& Service}

    \begin{twocolentry}{\textit{Fall 2024 -- Present}}
        \textbf{Tutor \& Grader}\\
        University of California San Diego
    \end{twocolentry}

    \vspace{0.10 cm}
    \begin{onecolentry}
        \begin{highlights}
            \item \textbf{Tutor:} DSC\,40A: Theoretical Foundations of Data Science I (Summer 2025)
            \item \textbf{Grader:} MATH\,180A/C, MATH\,181A, MATH\,185 (Fall 2024 -- Fall 2025, 7 quarters)
        \end{highlights}
    \end{onecolentry}

    %--------------------------------------------------------------------------
    %  Technical Skills Section
    %--------------------------------------------------------------------------
    \section{Technical Skills}

    \begin{onecolentry}
        \textbf{Languages:} Python, R, SQL, Java, JavaScript, HTML\\[2pt]
        \textbf{ML \& LLM:} PyTorch, PyTorch Geometric, Transformers, scikit-learn, LLM fine-tuning and post-training, RAG systems, memory-augmented LLMs, reasoning-based retrieval, continual learning\\[2pt]
        \textbf{Data \& Infrastructure:} Pandas, Spark, AWS, Docker, distributed training, vector databases
    \end{onecolentry}

\end{document}